% answers_short.tex —— 2005 年·普通高等学校招生考试·辽宁卷 · 简短答案速查
% 由 tools/make_short_answers.py 自 answers.tex 的 \gksolhead/\gkanswer 提取，
% 并与 2005/2005_数学真题_辽宁_答案要点.md 逐题比对后生成，勿手改（改源 MD 后重跑）。

\begin{gknotice}
\textbf{使用说明}：本节只给最终答案便于快速核对，完整推导见解析版。答案为编写者独立推导，\textbf{非官方标准答案}。
\end{gknotice}

\gkpart{选择题}
\begin{center}\begin{tabular}{ccccccccccccc}
\toprule
题号 & 1 & 2 & 3 & 4 & 5 & 6 & 7 & 8 & 9 & 10 & 11 & 12 \\
\midrule
答案 & B & B & D & D & C & C & C & B & A & A & B & A \\
\bottomrule
\end{tabular}\end{center}

\gkpart{填空题}
\begin{itemize}[leftmargin=2.6em]
  \item[\textbf{13.}] $n$ 为偶数时常数项为 $\mathrm{C}_{n}^{\frac{n}{2}}(-2)^{\frac{n}{2}}$；$n$ 为奇数时展开式中不存在常数项。
  \item[\textbf{14.}] $\dfrac{2}{3}$。
  \item[\textbf{15.}] $576$。
  \item[\textbf{16.}] $(\pi, 2\pi]$。
\end{itemize}

\gkpart{解答题}
\begin{itemize}[leftmargin=2.6em]
  \item[\textbf{17.}] (1) 由 $\overrightarrow{PC}\cdot\overrightarrow{PA}=0$ 得 $PC \perp PA$，又 $PC \perp AB$，故 $PC \perp$ 平面 $PAB$；(2) $\dfrac{\sqrt{3}}{3}$；(3) $2\sqrt{2}$
  \item[\textbf{18.}] (1) $S(\theta)=\sin 2\theta - \sin^{2}\theta$，$\theta \in \left(0, \dfrac{\pi}{4}\right)$；(2) 当 $\tan\theta=\dfrac{\sqrt{5}-1}{2}$ 时 $S$ 取最大值 $\dfrac{\sqrt{5}-1}{2}$
  \item[\textbf{19.}] (1) 由 $\dfrac{b_{n+1}}{b_n}=\dfrac{\sqrt{3}-1}{a_n+1}<\dfrac{\sqrt{3}-1}{2}$ 及 $b_1=\sqrt{3}-1$ 用数学归纳法得 $b_n \leqslant \dfrac{(\sqrt{3}-1)^{n}}{2^{n-1}}$；(2) 对 (1) 的界按等比级数求和得 $S_n<\dfrac{2\sqrt{3}}{3}$
  \item[\textbf{20.}] (1) $P_\text{\text{甲}}=0.68$，$P_\text{\text{乙}}=0.6$；(2) $\xi$ 的分布列为 $P(\xi=5)=0.68$、$P(\xi=2.5)=0.32$，$E\xi=4.2$；$\eta$ 的分布列为 $P(\eta=2.5)=0.6$、$P(\eta=1.5)=0.4$，$E\eta=2.1$（单位：万元）；(3) 当 $x=5$、$y=0$ 或 $x=4$、$y=2$ 时 $z$ 取最大值 $21$ 万元
  \item[\textbf{21.}] (1) 由椭圆方程与 $b^2=a^2-c^2$ 化简 $\left|\overrightarrow{F_1P}\right|=\sqrt{(x+c)^2+y^2}$ 得 $\left|\overrightarrow{F_1P}\right|=a+\dfrac{c}{a}x$；(2) $x^2+y^2=a^2$；(3) 当 $b^2 \leqslant ac$ 时存在这样的点 $M$，且 $\tan\angle F_1MF_2=2$；当 $b^2>ac$ 时不存在
  \item[\textbf{22.}] (1) $m=f(x_0)-x_0f'(x_0)$；(2) 令 $F(x)=g(x)-f(x)$，由 $f'$ 递减得 $F$ 在 $x_0$ 处取最小值 $0$，故 $g(x) \geqslant f(x)$；(3) $b \in \left[\dfrac{2-\sqrt{2}}{4}, \dfrac{2+\sqrt{2}}{4}\right]$，且 $a$、$b$ 满足 $b \geqslant \dfrac{1}{2a^{2}}$、$b \leqslant 1-\dfrac{a^{2}}{4}$
\end{itemize}
